\chapter{}

Ghìsola aveva ripreso la sua strada verso il campo, con un'ebrezza che empiva di gioia tutto il suo essere. Il movimento delle gambe assecondava questa ebrezza; e le sottane erano così lievi che non le sentiva nè meno.

Ella non si fidava d'Antonio che era capace di ridire tutto al padrone; non faceva nessun conto di Pietro; ed Agostino le piaceva più di tutti e tre.

In quel mentre questi, correndo attraverso i filari delle viti, e saltando le passate del grano nuovo, le andò incontro come quando con un palo in mano sfondava le zucche. Era in maniche di camicia, con i polsi tondi e forti e le vene strette dalla carne soda. Non portava il cappello; e gli occhi verdognoli, di una lucentezza di diaccio, sembravano senza palpebre.

Le saltò addosso e la gettò a terra; facendola piangere. Allora le chiese, per celia:

— Hai sentito male?

— Niente! Niente!

E lesta, alzandosi, lo afferrò a mezza vita; per fare altrettanto a lui. Ma Agostino le tirò giù le braccia. Ella sorrise, con il viso bagnato di lagrime; volle svignarsela; e puntò i piedi serrandoli insieme. Sicuro della sua forza, il giovine le gridava dentro gli orecchi:

— Ti faccio quello che voglio io! Non ruzzo. Tu lo sai!

Ella, allora, gli azzannò un braccio. Agostino, spingendo il braccio, le piegò la testa indietro, costringendola ad aprire i denti. Poi, piuttosto in collera, le domandò:

— Ed ora che cosa fai?

Ghìsola rispose, dopo aver sputato:

— Son la più debole. Te ne vanti? Com’è salata la tua pelle!

Egli la guardò negli occhi, per impaurirla.

— Quant’è che non vedi Pietro?

Ella cavò fuori la punta della lingua.

— Non viene più!

Egli che, da casa, lo aveva riconosciuto al vestito, ed era venuto per vederlo, le rifece la voce:

— Da vero?

— È quanto mi pare!

— Credevi che volessi venire a mangiare le ciliegie con lui?

E le andò addossò un’altra volta, per pestarle la punta delle scarpe tutte rotte lungo le ricuciture.

— Perchè non mi hai detto la verità? Con gli altri devi esser bugiarda; con me no.

E seguitava a farla indietreggiare. Ma ambedue caddero; battendo la fronte insieme. Allora egli ebbe il desiderio di litigare da vero; ma udì la sonagliera della sua mula:

— È il mio fratello che torna!

Si drizzò in ginocchio, per ascoltare meglio. Poi finì d'alzarsi e se ne andò vociando:

— Se l'ha strapazzata troppo!... Se l'ha strapazzata troppo! Non la sa guidare.

Il ciuffo a punta de' suoi capelli sudati gli sbatteva su le ciglia; e, con quegli orecchi stretti, tutta la testa, rotonda di dietro, sembrava una palla.

Ghìsola era rimasta lì, pentita di trovarsi stesa in terra a quel modo. Si alzò in fretta, pulendosi e guardandosi i polpacci delle mani chiuse a pugno; come quando era a tagliare l'erba e si riposava.

Quando era a tagliare l'erba ficcava la punta del falcino nel tronco di un albero, assettandosi un poco le vesti addosso, specie la camicietta che si sbottonava sempre; stringendo tra i denti le forcelle che una per volta ripigliava per mettersele nei capelli unti d'olio. Dopo aver toccato la punta del falcino, umida del legno lacerato, come di una saliva, cominciava a cantare; interrompendosi, e stando dritta in piedi. Poi, si sputava nelle mani e si rimetteva giù.

Talvolta, le veniva voglia di nascondere tutto il viso; e di restare così; di non esser veduta che dall'aria; di non mangiare più, di morire senza accorgersene.

Le veniva anche voglia di gridare; e aveva

paura.
